\section{The geometry of adding a site}
\label{sec:geometry}

This section describes what identification looks like in the space of
site profiles. Every identification statement in the
selection problem reduces to a statement about the affine hull
\begin{equation}
  \affhull(\pool)
  =\Bigl\{\textstyle\sum_{s=1}^S w_s c_s:\ \sum_{s=1}^S w_s=1\Bigr\},
  \label{eq:affine-hull}
\end{equation}
the set of affine combinations of the observed profiles. Its direction
space is exactly $\Vspan$, the span of the centered profiles from
\eqref{eq:span}.

The first step relates the target loading to the hull. The loading
$\ellvec=\ct-\E_{\tP}[\muC(X)]$ subtracts from the target profile a
point that is always a mixture of observed profiles, because
$\muC(x)=\sum_s \pr(C=c_s\mid X=x)\,c_s$ is a convex combination of them
at every $x$. The unidentified part of the loading is therefore the part of the target
profile itself that lies outside the hull.

\begin{lemma}[The loading and the affine hull]
\label{lem:loading}
Under \cref{ass:model} and the eigenvalue condition of
\cref{sec:setting}, $\ellvec\in\range(\CovC)$ if and only if
$\ct\in\affhull(\pool)$. Moreover, for the orthogonal projection
$\Pi^\perp$ onto $\Vspan^\perp$,
\begin{equation}
  \Pi^\perp\ellvec=\Pi^\perp(\ct-c_1),
  \qquad
  \norm{\Pi^\perp\ellvec}_2=\dist\bigl(\ct,\affhull(\pool)\bigr).
  \label{eq:loading-distance}
\end{equation}
\end{lemma}

\noindent
See \cref{app:proofs} for all proofs. \cref{lem:loading} turns the point-identification condition of the
companion paper into a statement an investigator can check from the
profiles alone. The pooled experiments determine the target effect
exactly when the target's context is an affine combination of the
observed contexts. The relevant set is the affine hull, not the convex
hull. A target outside the convex hull but inside the affine hull is
still point identified, because the fixed effects are linear in context,
so the model allows linear extrapolation along observed directions.
Identification fails only in a direction that no combination of observed
sites reaches.

The same projection governs the width under any restriction. On the
feasible set the projection of the treatment shift onto $\Vspan$ is
fixed, so only the part of the loading outside $\Vspan$ can move the
objective.

\begin{corollary}[The width depends on the target only through its offset]
\label{cor:width-offset}
Under the conditions of \cref{lem:loading} and for any restriction set
$\restrset$, the width $\width$ of \eqref{eq:endpoint-programs} depends
on the target $(\tP,\ct)$ only through $\Pi^\perp(\ct-c_1)$. In
particular it does not depend on the target covariate law $\tP$.
\end{corollary}

\noindent
The endpoints themselves do depend on $\tP$ through $\bT$ and the
in-span part of $\ellvec$, but both endpoints move by the same amount.
\cref{cor:width-offset} is what justifies writing the width as
$\width(\ct;\pool)$ and stating the target region below as a set of
profiles.

Without restrictions, partial identification in the response model
reduces to a dichotomy. \citet{boyarsky2026} formalize this.

\begin{proposition}[Unrestricted dichotomy]
\label{prop:dichotomy}
Let $\restrset=\R^{2p}$. Under the conditions of \cref{lem:loading},
\begin{equation}
  \width(\ct;\pool)=
  \begin{cases}
    0 & \text{if } \ct\in\affhull(\pool),\\
    +\infty & \text{otherwise.}
  \end{cases}
  \label{eq:dichotomy}
\end{equation}
\end{proposition}

\noindent
With $\restrset=\R^{2p}$ the program is always feasible under the model,
since $\beta^\circ$ satisfies the moment equations, so the width is never
undefined.

A candidate site enters this geometry through one object only, the
direction its profile adds to the span.

\begin{proposition}[Span expansion is affine-hull expansion]
\label{prop:hull}
Let $\newdir{k}$ be the component of $\cand{k}-c_1$ orthogonal to
$\Vspan$. Then $\Vspan'$ of \eqref{eq:span-update} satisfies
$\dim\Vspan'=r+1$ if and only if $\newdir{k}\neq0$, which holds if and
only if $\cand{k}\notin\affhull(\pool)$. In that case
$\Vspan'^{\perp}=\Vspan^\perp\cap\{\newdir{k}\}^\perp$, and under the
augmented eigenvalue condition of \cref{subsec:augmentation} the
augmented covariance satisfies $\range(\CovC')=\Vspan'$ and
$\kernel(\CovC')=\Vspan'^\perp$. If instead
$\cand{k}\in\affhull(\pool)$, the span, the range, and the kernel are
unchanged.
\end{proposition}

Together with \cref{prop:dichotomy}, this gives the unrestricted
selection rule. A candidate collapses an infinite width at target $\ct$
to zero exactly when the target's offset from the hull,
$\Pi^\perp(\ct-c_1)$, is a nonzero multiple of $\newdir{k}$. When the
offsets of the targets span two or more directions outside $\Vspan$, no
single candidate gives point identification, and every candidate leaves
the unrestricted minimax width at $+\infty$. The unrestricted rule then
cannot rank candidates, which is one reason the selection problem needs
restrictions. Only the direction of $\newdir{k}$ enters. How far outside the hull the candidate sits does
not change what becomes identified, because a fixed-effects model pins
the same coefficient combination from a small step along a direction as
from a large one. Distance along the new direction affects the precision of the added
constraint, through the context variance the new site contributes. A
farther site therefore improves precision but not identification.

Restrictions replace the dichotomy with finite widths. In that case the
distance from the target to the hull is the right continuous measure of
what remains unidentified. The simplest statement uses a cap on the size
of the treatment shift $\shift=\beta_1-\beta_0$.

\begin{proposition}[Width equals distance under a norm cap]
\label{prop:distance}
Let the restriction be $\norm{\shift}_2\leq\tau$ with no restriction on
$\beta_0$, let $\bar\shift$ be the minimum-norm solution of
$\CovC\shift=\dmom{1}-\dmom{0}$, and assume
$\tau\geq\norm{\bar\shift}_2$. Under the conditions of
\cref{lem:loading},
\begin{equation}
  \width(\ct;\pool)
  =2\sqrt{\tau^2-\norm{\bar\shift}_2^2}\;
  \dist\bigl(\ct,\affhull(\pool)\bigr).
  \label{eq:width-distance}
\end{equation}
\end{proposition}

\noindent
The same formula holds for the augmented pool, with two changes. The
distance is to $\affhull(\pool\cup\{k\})$, and the minimum-norm shift
becomes $\Pi_{\Vspan'}\shift^\circ$, where $\shift^\circ$ is the true
shift. Writing $\newdir{k}$ with unit length,
\begin{equation}
  \width(\ct;\pool\cup\{k\})
  =2\sqrt{\tau^2-\norm{\bar\shift}_2^2-(\newdir{k}^\top\shift^\circ)^2}\;
  \dist\bigl(\ct,\affhull(\pool\cup\{k\})\bigr),
  \label{eq:width-distance-aug}
\end{equation}
provided the true shift satisfies the cap (\cref{app:proof-distance}).
The first factor depends on the candidate through the unknown scalar
$\newdir{k}^\top\shift^\circ=\reveal{1}{k}-\reveal{0}{k}$ that its
experiment reveals. So under this restriction the ranking of candidates
by \eqref{eq:minimax} is not a ranking by distance alone. Two statements
do hold before the experiment. First, at the neutral forecast of
\cref{sec:selection}, which sets this scalar to zero, the first factor is
the same for every candidate, and the minimax rule chooses the candidate
whose profile, added to the hull, leaves no target in $\targset$ far
from it. Second, because the first factor is largest at zero, the
width at the neutral forecast is an upper bound on each candidate's
realized width under this restriction.
